\Chapter{40}

\Verse{1}
{
	\hP{}
}
{
	\eP{
		40

¹ And YHWH spoke to Moses, saying,\\
	}
}
{
}

\Verse{2}
{
	\hP{}
}
{
	\eP{
		“On the day of the first month, on the first of the month, you shall
set up the Tabernacle of the Tent of Meeting.\\
	}
}
{
}

\Verse{3}
{
	\hP{}
}
{
	\eP{
		And you shall set the Ark of the Testimony there and have the pavilion
cover over the ark.\\
	}
}
{
}

\Verse{4}
{
	\hP{}
}
{
	\eP{
		And you shall bring the table and make its arrangement, and bring the
menorah and put up its lamps.\\
	}
}
{
}

\Verse{5}
{
	\hP{}
}
{
	\eP{
		And you shall put the golden altar for incense in front of the Ark of
the Testimony, and you shall set the cover of the entrance of the
Tabernacle.\\
	}
}
{
}

\Verse{6}
{
	\hP{}
}
{
	\eP{
		And you shall put the altar of burnt offering in front of the entrance
of the Tabernacle of the Tent of Meeting.\\
	}
}
{
}

\Verse{7}
{
	\hP{}
}
{
	\eP{
		And you shall put the basin between the Tent of Meeting and the altar
and put water there.\\
	}
}
{
}

\Verse{8}
{
	\hP{}
}
{
	\eP{
		And you shall set the courtyard all around and put on the cover of the
courtyard’s gate.\\
	}
}
{
}

\Verse{9}
{
	\hP{}
}
{
	\eP{
		And you shall take the anointing oil and anoint the Tabernacle and
everything that is in it, and you shall make it and all of its equipment
holy, and it will be holiness.\\
	}
}
{
}

\Verse{10}
{
	\hP{}
}
{
	\eP{
		And you shall anoint the altar of burnt offering and all of its
equipment, and you shall make the altar holy: and the altar will be holy
of holies.\\
	}
}
{
}

\Verse{11}
{
	\hP{}
}
{
	\eP{
		And you shall anoint the basin and its stand and make it holy.\\
	}
}
{
}

\Verse{12}
{
	\hP{}
}
{
	\eP{
		“And you shall bring Aaron and his sons forward to the entrance of the
Tent of Meeting and wash them with water.\\
	}
}
{
}

\Verse{13}
{
	\hP{}
}
{
	\eP{
		And you shall dress Aaron with the holy clothes and anoint him and
make him holy, and he shall function as a priest for me.\\
	}
}
{
}

\Verse{14}
{
	\hP{}
}
{
	\eP{
		And you shall bring his sons forward and dress them with coats\\
	}
}
{
}

\Verse{15}
{
	\hP{}
}
{
	\eP{
		and anoint them as you anointed their father, and they shall function
as priests for me. And it will be for their anointing to be theirs as an
eternal priesthood through their generations.”\\
	}
}
{
}

\Verse{16}
{
	\hP{}
}
{
	\eP{
		And Moses did it. According to everything that YHWH commanded him, he
did so.\\
	}
}
{
}

\Verse{17}
{
	\hP{}
}
{
	\eP{
		And it was: in the first month, in the second year, on the first of
the month, the Tabernacle was set up.\\
	}
}
{
}

\Verse{18}
{
	\hP{}
}
{
	\eP{
		And Moses set up the Tabernacle and put on its bases and set its
frames and put on its bars and set up its columns.\\
	}
}
{
}

\Verse{19}
{
	\hP{}
}
{
	\eP{
		And he spread the Tent over the Tabernacle and set the Tent’s covering
on it above, as YHWH had commanded Moses.\\
	}
}
{
}

\Verse{20}
{
	\hP{}
}
{
	\eP{
		And he took the Testimony and put it into the ark, and he set the
poles on the ark, and he put the atonement dais on the ark above.\\
	}
}
{
}

\Verse{21}
{
	\hP{}
}
{
	\eP{
		And he brought the ark into the Tabernacle and set the covering
pavilion and covered over the Ark of the Testimony, as YHWH had
commanded Moses.\\
	}
}
{
}

\Verse{22}
{
	\hP{}
}
{
	\eP{
		And he put the table in the Tent of Meeting on the northward side of
the Tabernacle, outside of the pavilion.\\
	}
}
{
}

\Verse{23}
{
	\hP{}
}
{
	\eP{
		And he made the arrangement of bread on it in front of YHWH, as YHWH
had commanded Moses.\\
	}
}
{
}

\Verse{24}
{
	\hP{}
}
{
	\eP{
		And he set the menorah in the Tent of Meeting opposite the table, on
the southward side of the Tabernacle,\\
	}
}
{
}

\Verse{25}
{
	\hP{}
}
{
	\eP{
		and he put up the lamps in front of YHWH, as YHWH had commanded Moses.\\
	}
}
{
}

\Verse{26}
{
	\hP{}
}
{
	\eP{
		And he set the golden altar in the Tent of Meeting in front of the
pavilion,\\
	}
}
{
}

\Verse{27}
{
	\hP{}
}
{
	\eP{
		and he burned the incense of fragrances on it, as YHWH had commanded
Moses.\\
	}
}
{
}

\Verse{28}
{
	\hP{}
}
{
	\eP{
		And he set the cover of the entrance of the Tabernacle.\\
	}
}
{
}

\Verse{29}
{
	\hP{}
}
{
	\eP{
		And he set the altar of burnt offering at the entrance of Tabernacle
of the Tent of Meeting, and he offered up the burnt offering and the
grain offering on it, as YHWH had commanded Moses.\\
	}
}
{
}

\Verse{30}
{
	\hP{}
}
{
	\eP{
		And he set the basin between the Tent of Meeting and the altar, and he
put water for washing there.\\
	}
}
{
}

\Verse{31}
{
	\hP{}
}
{
	\eP{
		And Moses and Aaron and his sons washed their hands and their feet
from it.\\
	}
}
{
}

\Verse{32}
{
	\hP{}
}
{
	\eP{
		When they came to the Tent of Meeting and when they came forward to
the altar they would wash, as YHWH had commanded Moses.\\
	}
}
{
}

\Verse{33}
{
	\hP{}
}
{
	\eP{
		And he set up the courtyard all around the Tabernacle and the altar,
and he put on the cover of the courtyard’s gate.

And Moses finished the work.\\
	}
}
{
}

\Verse{34}
{
	\hP{}
}
{
	\eP{
		And the cloud covered the Tent of Meeting, and YHWH’s glory filled the
Tabernacle.\\
	}
}
{
}

\Verse{35}
{
	\hP{}
}
{
	\eP{
		And Moses was not able to come into the Tent of Meeting, because the
cloud had settled on it and YHWH’s glory filled the Tabernacle.\\
	}
}
{
}

\Verse{36}
{
	\hP{}
}
{
	\eP{
		And when the cloud was lifted from on the Tabernacle, the children of
Israel would travel—in all their travels—\\
	}
}
{
}

\Verse{37}
{
	\hP{}
}
{
	\eP{
		and if the cloud would not be lifted, then they would not travel until
the day that it would be lifted.\\
	}
}
{
}

\Verse{38}
{
	\hP{}
}
{
	\eP{
		Because YHWH’s cloud was on the Tabernacle by day, and fire would be
in it at night, before the eyes of all the house of Israel in all their
travels.

------------------------------------------------------------------------

COMMENTARY ON THE TORAH

------------------------------------------------------------------------

After the rush of forces in Exodus, Leviticus is a rather tranquil book.
Following the mighty acts of God in Exodus, Leviticus is
disproportionately the words of God. There is very little
narrative—barely three chapters out of twenty-seven in the book—and no
poetry (with the possible exception of a few isolated verses). More than
the other four of the Five Books of Moses, Leviticus pictures the deity
speaking. Consistent with this is the absence of movement in Leviticus.
Unlike Genesis and Exodus, the entire book of Leviticus is set in one
place: at the foot of Mount Sinai. In Exodus, the universe—the land,
waters, and sky—is in disarray. Following that, Leviticus is concerned
with orderliness, arrangement.

Leviticus pictures the beginning of the solidification—the formation of
an identity—of a people through law. The law is significant as a subject
of study in itself and in context of the study of history as well. But,
also in narrative terms, it is important to see the place of the law in
the development of the story of the people: laws of how to behave toward
other human beings, laws that bring identity (what to eat, religious
ceremonies, holidays). The laws are not merely listed. They are narrated
as words of YHWH to Moses (and Aaron) and thus are presented as an
integrated component of the account.

In Leviticus the knowledge of God, which is an explicit, continuing
component in Exodus, likewise resides in the background. The term “to
know” (Hebrew yd‘) never occurs in Leviticus in the sense of knowledge
of God. YHWH now is known to Israel. This is assumed.

Impressively, the story stops being about Moses in Leviticus. The laws
are conveyed through him, certainly, but the book is about the laws
themselves and the fact of their being revealed to the people. The only
stories are primarily about Aaron and his sons. More broadly, Leviticus
does not develop personalities. It is about institutions, not
individuals. And this, too, contributes to the reduction of tension in
Leviticus following the great confrontations of Exodus.

All of this contributes to the change of feeling of this book. Leviticus
is both integral to the context of the books that precede and follow it
and, at the same time, anomalous in character.

------------------------------------------------------------------------

Leviticus

------------------------------------------------------------------------

Unlike in Exodus, there are few miracles portrayed in Leviticus, yet the
book, rich in background like the two books that precede it, is pervaded
with the presence of the divine. The continuing scene has already been
set in Exodus: A column of cloud or fire is visible at all times. A Tent
of Meeting, the Tabernacle, stands outside the camp. Moses goes in and
out of the tent, his face veiled and then unveiled, bringing laws from
God. The book begins with God speaking from the Tabernacle, thus
establishing the Tabernacle from the outset as the visible, present
entity among the people through which the divine word enters the world.
The cosmic/miraculous thus looms, and a channel between it and the
earthly operates.\\
	}
}
{
}

